% !TEX root = ../main.tex

Il file \textit{io\_utils.py} raccoglie le funzioni di caricamento dati. Espone un'unica funzione, \texttt{load\_csv}, che avvolge \texttt{pd.read\_csv} aggiungendo validazione delle colonne e gestione dei valori mancanti.

\section{\texorpdfstring{load\_csv(path, \ldots)}{load\_csv(path, ...)}}
\label{sec:load-csv}
\begin{apidesc}
  \item[\textbf{Scopo}]
    Caricare un file CSV in un \texttt{DataFrame}, con supporto a separatori e decimali personalizzati, rinomina delle colonne e controllo dei valori mancanti.
  \item[\textbf{Parametri}]\hfill
    \begin{description}[font=\normalfont\ttfamily\bfseries, labelwidth=1.5em, leftmargin=2em, itemsep=1pt, parsep=0pt, topsep=1pt]
      \item[path] str o Path $\to$ percorso del file CSV
      \item[sep] str, default \texttt{","} $\to$ separatore di colonna
      \item[decimal] str, default \texttt{"."} $\to$ separatore decimale
      \item[rename\_columns] dict o \texttt{None}, default \texttt{None} $\to$ mappa \texttt{\{vecchio: nuovo\}} per rinominare le colonne
      \item[required\_columns] list o \texttt{None}, default \texttt{None} $\to$ lista di colonne obbligatorie
      \item[drop\_missing] bool, default \texttt{False} $\to$ se \texttt{True} elimina le righe con NaN; altrimenti solleva eccezione
    \end{description}
  \item[\textbf{Note}]
    Le righe che iniziano con \texttt{\#} vengono ignorate; gli spazi dopo il separatore vengono rimossi automaticamente. Questi comportamenti sono fissi.
  \item[\textbf{Restituisce}]
    \texttt{pandas.DataFrame} con i dati del file.
  \item[\textbf{Eccezioni}]
    \texttt{ValueError} se una o più colonne di \texttt{required\_columns} sono assenti;
    \texttt{ValueError} se il file contiene valori NaN e \texttt{drop\_missing=False}.
\end{apidesc}
\vspace{-15pt}
\paragraph{Implementazione}
\begin{minted}{python}
def load_csv(path, sep=",", decimal=".", rename_columns=None,
             required_columns=None, drop_missing=False):
    df = pd.read_csv(Path(path), sep=sep, decimal=decimal,
                     comment='#', skipinitialspace=True)
    if rename_columns:
        df = df.rename(columns=rename_columns)
    if required_columns:
        missing = [col for col in required_columns if col not in df.columns]
        if missing:
            raise ValueError(f"Colonne mancanti: {missing}")
    if drop_missing:
        df = df.dropna()
    elif df.isna().any().any():
        raise ValueError("Il file contiene valori mancanti")
    return df
\end{minted}
\vspace{-10pt}
\paragraph{Esempio d'uso}
\begin{minted}{python}
# CSV con separatore ";" e decimale "," (formato italiano)
df = load_csv("misure.csv", sep=";", decimal=",",
              rename_columns={"t [s]": "t", "x [m]": "x"},
              required_columns=["t", "x"], drop_missing=True)
\end{minted}
